problem:  
Let \(M^n \subset S^{n+p}(1)\) be a compact minimal submanifold, with \(\|B\|^2\) the squared norm of its second fundamental form and \(\lambda_1(M)\) the first nonzero Laplace–Beltrami eigenvalue on \(M\). By the classical result of Takahashi, the coordinate functions restricted to \(M\) are eigenfunctions with eigenvalue \(n\), hence \(\lambda_1(M) \le n\).  

It is conjectured (Yau, 1980s) that for any compact minimal hypersurface \(M^n \subset S^{n+1}(1)\), equality \(\lambda_1(M) = n\) actually holds; partial results and examples support this conjecture. On the other hand, Simons’ theorem and its refinements show that when \(\|B\|^2\) is “small,” rigidity and gap phenomena occur.  

**Problem (Coupled Eigenvalue–Curvature Gap Conjecture):**  
For each pair \((n,p)\), is there a positive constant \(\varepsilon_{n,p}\) such that the following holds?  
If \(M^n \subset S^{n+p}(1)\) is a compact minimal submanifold satisfying  
\[
n - \lambda_1(M) < \varepsilon_{n,p}
\quad\text{and}\quad
\|B\|^2 < n + \varepsilon_{n,p},
\]
then \(M\) must be congruent to one of the known *homogeneous minimal submanifolds* (such as Clifford minimal hypersurfaces or the Veronese surface).  

Equivalently, can simultaneous near-equality in both the intrinsic first eigenvalue and the extrinsic curvature functional force geometric homogeneity?  

Why is it a "good" problem:  
This formulation corrects the spectral facts and makes the conjecture conceptually and mathematically sound. It is a “good” problem because:

1. **Bridges two fundamental rigidity principles:**  
   It combines Simons-type extrinsic curvature gaps with Yau’s intrinsic eigenvalue rigidity conjecture. Establishing or disproving a *finite double pinching domain* would reveal how intrinsic and extrinsic data jointly constrain the geometry of minimal immersions in spheres.

2. **Conceptually simple, yet deeply open:**  
   The conditions involve only two classical invariants (\(\lambda_1\) and \(\|B\|^2\)), but proving any nontrivial interaction requires deep analytic and geometric tools—connecting Simons’ identity, Reilly-type inequalities, and possibly stability operator spectra.

3. **Anchored in known structures but extends them radically:**  
   Each component (the eigenvalue equality \(\lambda_1=n\), the Simons gap \(\|B\|^2=n\)) is well studied separately; their intersection has not been explored systematically. This new two-parameter gap formulation extends current theory in a natural yet uncharted direction.

4. **Potential for field unification:**  
   A resolution would link spectral geometry, minimal submanifold theory, and global analysis, pioneering new rigidity mechanisms that might apply beyond spheres, for example to Einstein manifolds or manifolds with positive Ricci curvature.

In essence, the problem asks whether *simultaneous analytic and geometric near-rigidity implies global symmetry*, an idea both elegant and formidable—an archetypal feature of a “good” research-level problem.